\paragraph{Following a bilingual request for historical evidence.}
A free modern writer drew four glyphs on each of 9,600 one- or two-form questions and answered none exactly. We paired each question with a four-row plate of attested historical forms. A trained source-addressing model answered 6,884 questions exactly. Its first image was correct on 4,777 of 4,800 single-form questions, but it added a second form 2,577 times. Learning which stage names occurred in the question raised complete answers to 9,563. A second question frame yielded 8,993, including 1,987 of 2,400 Chinese single-form requests.

Supervising the extent of a pictured stage name yielded 9,565 exact answers (99.64\%) in each of three frames on the same 600 plates. All 249 development families and 1,372 distinct historical bitmaps were held out from training. The model answered both orders correctly in 2,388 of 2,400 pairs. Moving all four ink payloads between the rows, retaining the headings and scoring the relocated pictures, preserved 9,565 correct responses. Erasing the question or the plate gave zero exact answers. All 35 errors involved a competing Chinese headword that also appeared in the stage name. The model retained the original stop and every answer pixel on all 9,135 fresh modern passages, including all 163 exact expressions. Figure~\ref{fig:historical-plates} shows the observed distinction between selecting visible historical evidence and independently composing a missing expression.
